\section{Buffers, Offsets, and Truncation}
\label{app:truncation}

The finite refinement keeps the buffer and the state at its start.
We prove that this is enough to replace the stems, then show how
to take a finite hole from the resulting periodic tuple.

\subsection{Safety after a buffer}

Along a ray with onset $a$ and period $p$, a cut at or after
$a$ is described by its reader state and its phase modulo $p$.
The next pair is determined by the current one. There are $Np$
possible pairs. Thus every pair in the tail occurs among the first
$Np$ cuts of the tail, by following the deterministic sequence
up to its first repetition. Every cut has a representative
strictly below $a+Np$ with the same state and infinite rest.
Cuts before the onset represent themselves.

A word of the original stopped tuple cannot have a suffix of
length $(C+1)p$ with period $p$. The suffixes of that suffix of
lengths $p,2p,\dots,(C+1)p$ would all be prefixes of the suffix
itself, contradicting depth at most $C$. Each of these lengths is
positive. Since $B\ge (N+C+1)p$ for every period $p\le D$, any
such suffix of a shortened stopped word would have survived from
the original. So no shortened stopped word has that periodic suffix.

A comparison of a shortened stopped target with a shortened stopped
host or buffered ray prefix lifts by the refinement. For a proper
target cut, the matching word contains no terminal marker and so
the lifted comparison occurs in the decoded words. For a final
target cut, use equality of endpoints and the visit property of
the refinement. If the lifted host cut is after its terminal
marker, the state is the same just before it. The empty rest of
the target fits there. Thus safety against stopped targets holds
on the shortened finite words and buffered prefixes. The stopped
words remain eligible, because endpoints and cut states lift.

Consider a match of any stopped target rest in a new ray. Move
the host cut back to its representative below the new onset
plus $Np$. If the match reaches beyond the buffer, its last
$(C+1)p$ letters lie in the periodic part. Indeed the host cut
is before the onset plus $Np$ and the buffer has at least
$Np+(C+1)p$ letters. So there are more than $(C+1)p$ letters in
the match, and those final letters start at or after the onset.
They form the excluded periodic suffix of the stopped word.
Every match at the representative thus fits in the buffered
prefix, with the same target and pair of states. This proves safety
of every prefix of the new rays against the stopped targets.

\subsection{Replacing the stems}

For any periodic tail, a period that works from a later cut
also works from its onset. To verify this, suppose the ray is
periodic with period $q>0$ from the onset and period $p>0$ from
a later cut. Shift two positions whose difference is $p$ forward
by the same multiple of $q$ until both are beyond that later cut.
Their letters agree there by period $p$, and shifting back
preserves them by period $q$. Thus we may use the loop bound
of \cref{lem:periodic} at the least onset. If the onset is positive,
the preceding letter differs from its copy one period later.
Otherwise that period would already hold one cut earlier.

For coordinate $i$ let $A_i$ be the old onset and $A_i'$ the
length of the new stem. It is followed by the same periodic text
as after $A_i$ of the old ray. If the original stem is empty,
so is the new one. If it is nonempty, its retained last letter
still breaks the period one position before the end of the new
stem. The period-extension argument therefore excludes every
earlier onset. Thus the onset of the new ray is $A_i'$.
The kept state at the end of the stem gives the same run in
the periodic text. Every cut before that state lifts by
refinement to a live cut in the old buffered word. The new
ray is live everywhere and eventually rejected, because the old
ray is rejected from some cut onward.

With old offsets $o_i$, use new offsets $o_i'=o_i+A_i'-A_i$.
Suppose a host cut $a$ of new ray $i$ and a target cut $b$ of
new ray $j$ have equal infinite rests and their states are admitted.
Either both cuts are in the periodic parts, or both lie before
the onsets at the same distance. This is because the first cut
from which a rest is purely periodic is a property of the
infinite rest itself. In the first case, take the old cuts
$a'=a-A_i'+A_i$ and $b'=b-A_j'+A_j$. They have the same rests
and states.

In the second case, the two new finite rests through the terminal
markers agree, since the buffers have the same length.
They form a comparison with nonempty target rest. Lift it to old cuts
$a',b'$ of the marked words. Since the terminal marker occurs
only as the last letter, the lifted rests agree and both have that
marker at the same distance. The distances to the markers
may differ from the distances in the new words. The original
rays after their buffers agree, since their new rays have equal
infinite rests and agree past their onsets with the old tails.
Hence the lifted cuts have equal infinite rests on the old
rays, with the same admitted states. Equality of distances
before the old markers implies $b'-a'=A_j-A_i$.
Equality of distances before the new onsets implies
$b-a=A_j'-A_i'$. In either case
\[
   b'-a' = b-a + (A_i'-A_i)-(A_j'-A_j).
\]
The old inequality $o_i+b'-a'+1\le o_j$ therefore gives
$o_i'+b-a+1\le o_j'$. This reasoning applies with $i=j$ and
also when $a=b$. Together with safety and eligibility it proves
that the replacements form a periodic tuple, completing
\cref{lem:buffered}.

\subsection{A bounded solution of the strict constraints}

Suppose the onsets are bounded by $A$, the periods by $D$, and
the stopped word lengths by $J$. Write $H=A+ND+1$.
Each ray reaches rejection below $H$, by taking a representative
of any rejected cut. In particular the state at $H-1$ is
rejected. At any cut $b\ge H$, the state is a live successor
of a rejected state. Hence any admitted target cut of a ray
has $b<H$.

Use all ordered pairs of coordinates to form a graph with edge
weight $\min(H,o_j-o_i)$ from $i$ to $j$, extending the offsets
arbitrarily to stopped coordinates. The old offsets dominate
every edge difference, so each directed cycle has total weight
at most zero by telescoping. Let $u_j$ be the maximum weight of
a directed walk ending at $j$, allowing the empty walk at $j$.
Deleting a closed part never decreases weight, and preserves the
endpoints, so the maximum is attained on a simple path. Its
weight lies in $[0,kH]$, since there are at most $k$ vertices
and each step has weight at most $H$. Appending the edge $i\to j$
gives $u_j-u_i\ge\min(H,o_j-o_i)$. For every equal-rest
admission between moving cuts, $b-a+1\le H$ by the target bound,
and $b-a+1\le o_j-o_i$ by \eqref{eq:offset}. So this entire
system of strict inequalities is satisfied by the $u_i$.
In particular it includes self-comparisons and identical cuts.
At an identical cut an admission would contradict
\eqref{eq:offset} itself.

Take the truncation lengths
\[
                T_i=u_i+H+A+2D+J
\]
for the moving coordinates. Adding the common constant preserves
the required differences. Each truncated ray is live and its
endpoint is rejected. The stopped words are eligible and
admissions against stopped targets are excluded by the model.
If a ray is a target in an admitted comparison, its target cut
is before $H$. Its rest has length greater than $J$, so a
stopped host cannot contain it.

For two rays, a match of this target rest gives more than
$A+2D$ letters of agreement. We verify that this forces agreement
of their infinite rests. Beyond their common initial portion up
to the later onset as measured from the two cuts, both rests are
purely periodic, with periods $p,q\le D$. This initial portion
has length at most $A$. The periodic words agree for at least
$p+q$ letters. On the first $q$ positions of the second word,
shifting by $p$ gives equal letters, by agreement with the
first word. Period $q$ extends those equalities to all positions
of the second word. Both words now have period $p$ and their first
$p$ letters agree. This proves equality of the infinite rests.
The constraint for this pair makes $T_i-a<T_j-b$, contradicting
the prefix relation. Hence all comparisons are safe, and the
tuple is a hole. The moving lengths are at most $kH+H+A+2D+J$,
and this bound includes each stopped word length. Summing proves
\cref{lem:truncation}.
